\problemname{Äta choklad}

\noindent
Olle och Lisa har fått $N$ stycken chokladkakor i olika former av sina föräldrar, som de nu planerar att äta tillsammans (chokladkakorna, inte föräldrarna).
Varje chokladkaka är en rektangel som består av $R_i \times C_i$ rutor.
Chokladkakorna ligger på ett bord, och för att äta upp dem kommer de turas om att göra följande:
\begin{itemize}
  \item Välj en chokladkaka från bordet.
  \item Bryt den i två delar, så att båda är rektanglar med heltalssidor.
  \item Ät upp den mindre delen.
  \item Lägg tillbaka den andra delen på bordet.
\end{itemize}
Om chokladkakan de väljer har storlek $1 \times 1$ äter de upp den direkt, och lägger inte någonting på bordet.
De har redan bestämt att Lisa får börja.



Båda är väldigt glupska, och försöker maximera antalet chokladrutor de själva äter.
De är båda perfekta logiker som vet att den andra också spelar efter samma mål.  
Olle är dessutom väldigt mån om rättvisa, och vill gärna att de totala mängderna rutor choklad de båda äter ska vara så lika som möjligt.
Innan de börjar kan han välja att slänga högst en chokladkaka i soptunnan utan att Lisa märker någonting.



Låt $V_O$ och $V_L$ vara antalet rutor Olle respektive Lisa äter. 
Din uppgift är att bestämma vad den minimala skillnaden, alltså $|V_O - V_L|$, är,
efter att Olle slängt högst en chokladkaka.


\section*{Indata}
På den första raden står talet $N$ ($2 \leq N \leq 10^5$), antalet chokladkakor på bordet.
Därefter följer $N$ rader, där den $i$:te raden har talen $R_i$ och $C_i$ ($1 \leq R_i, C_i \leq 10^6$).


\section*{Utdata}
Skriv ut ett heltal, det minimala värdet på $|V_O - V_L|$ som beskrivet ovan.


\section*{Poängsättning}


\noindent
\begin{tabular}{| l | l | p{12cm} |}
  \hline
  \textbf{Grupp} & \textbf{Poäng} & \textbf{Gränser} \\ \hline
  $1$    & $10$        & Alla chokladkakor har samma storlek. \\ \hline
  $2$    & $20$       & $R_i, C_i \leq 10$ \\ \hline
  $3$    & $15$       & $N \leq 50$ \\ \hline
  $3$    & $15$       & $N \leq 500$ \\ \hline
  $4$    & $40$       & Inga ytterligare begränsningar. \\ \hline
\end{tabular}

\section*{Förklaring av exempelfall}
